% Sections 5.1 to 5.4, from lectures/L05/appendix/a_variables_and_hardware.md: one section per
% part, A.1 to A.4.

\section{What This Chapter Is About}
\label{c5:a1}
Every construct you need in order to describe hardware, you have already used: \code{entity} and
\code{architecture} in Chapter~\ref{c1:ch}, \code{std_logic_vector}, \code{process} and \code{case}
in Chapter~\ref{c2:ch}, the clocked process and the scheduling rule in Chapter~\ref{c3:ch},
subcomponents and generics in Chapters~\ref{c2:ch} and~\ref{c4:ch}.

Two things are left, and they are the two hardest to work out on your own:
\begin{itemize}
  \item The first is \code{variable}, the one construct the book has not needed until now:
    \begin{itemize}
      \item It looks like a signal and behaves nothing like one, and the difference is exactly the
        kind that produces a design which simulates plausibly and is quietly wrong.
      \item \Secref{c5:a2} works through the case where it bites.
    \end{itemize}
  \item The second is what any of this actually becomes:
    \begin{itemize}
      \item Since Chapter~\ref{c1:ch} you have taken it on trust that the synthesis tool builds the
        circuit you described. It does, but not out of anything you would recognize.
      \item \Secref{c5:a3} opens the box. An FPGA turns out to contain no gates at all, which is why
        Chapter~\ref{c2:ch}'s Karnaugh maps bought you less than they seemed to. Your clock has a
        speed limit, and something specific sets it. And a \code{case} with a branch missing will
        quietly hand you a memory element nobody asked for.
    \end{itemize}
\end{itemize}
Neither section answers ``is my logic right?'' You have been answering that since
Chapter~\ref{c1:ch}. They answer the two questions after it: is this what I meant, and will it fit
and run.

\section{\code{variable}: Process-Local Storage with an Immediate Update}
\label{c5:a2}
A \code{variable} is declared inside a \code{process}, between \code{process(...)} and
\code{begin}. The critical difference from a \code{signal}: variable assignment uses \code{:=} and
takes effect \textbf{immediately}, so the very next statement that reads it sees the updated value,
exactly like an assignment in C. There is no scheduling, none of the
wait-until-the-process-suspends delay that \secref{c3:a6} spelled out for signals, and no ``last
assignment wins''.

\begin{booktable}{@{}p{9.5em}LL@{}}
  & \thead{\code{signal} (\code{<=})} & \thead{\code{variable} (\code{:=})} \\
  \midrule
  Declared & Architecture's declarative part & Process's declarative part \\
  Visible from & Anywhere in the architecture & Only within its own process \\
  Update timing & Scheduled: applied when the process next suspends & Immediate: applied before
    the next statement \\
  Repeated assignment in one pass & Only the last one survives; earlier ones are discarded & Every
    one takes effect, in order \\
\end{booktable}

\subsection{Why this actually matters: an XOR-parity generator}
\label{c5:parity}
An entity computing the XOR-parity of an 8-bit input, \code{'1'} if \code{bits} has an odd number of
set bits:

\begin{vhdlcode}
entity parity_gen is
    port(bits  : in  std_logic_vector(7 downto 0);
         parity: out std_logic);
end entity;
\end{vhdlcode}

\noindent Coming from software, the first instinct is to accumulate the running XOR in a loop, using
a signal as the accumulator:

\begin{vhdlcode}
-- Don't do this: shown to demonstrate why it doesn't work.
architecture broken of parity_gen is
signal parity_s: std_logic;
begin
    parity <= parity_s;

    process(bits) is
    begin
        parity_s <= '0';
        for i in bits'range loop
            parity_s <= parity_s xor bits(i);
        end loop;
    end process;
end architecture;
\end{vhdlcode}

\noindent Trace one pass, using the scheduling rule from \secref{c3:a6}: every read of
\code{parity_s} returns whatever it held \emph{before this pass started} (call it
\code{old_parity}), and only the last assignment survives.
\begin{itemize}
  \item \code{parity_s <= '0';} schedules \code{'0'}. Not yet applied.
  \item \code{i = 7} (\code{bits'range} on \code{(7 downto 0)} runs high-to-low): reads
    \code{parity_s}, gets \code{old_parity}, because the \code{'0'} has not applied. Schedules
    \code{old_parity xor bits(7)}, discarding the \code{'0'}.
  \item \code{i = 6}: reads \code{parity_s}, again \code{old_parity}. Schedules
    \code{old_parity xor bits(6)}, discarding the previous one.
  \item \ldots and so on for every iteration \ldots
  \item \code{i = 0}, the last statement in the pass: still reads \code{old_parity}, schedules
    \code{old_parity xor bits(0)}, discarding everything before it.
  \item The process suspends and \code{parity_s} becomes \code{old_parity xor bits(0)}, the only one
    of nine scheduled assignments that survived.
\end{itemize}
The result depends only on \code{bits(0)} and on whatever \code{parity_s} happened to already hold.
Seven iterations and the initial \code{'0'} were computed and thrown away. This is not a corner
case: it is what a signal-as-accumulator \emph{always} does, because every iteration reads the same
frozen value and only the textually last assignment is kept.

There is a second thing wrong with it, which the trace does not show. The surviving assignment is
\code{parity_s <= parity_s xor bits(0);}, and \code{parity_s} is not in the sensitivity list: as
hardware that is a \textbf{combinational feedback loop}, an output fed back to its own input through
a gate with no flip-flop in the path. GHDL analyzes and simulates it happily, which is how the bug
reaches a simulation at all; Quartus reports the loop rather than building it. The lesson is the
same one either way, and the fix below removes both problems at once.

The same computation with a \code{variable}:

\begin{vhdlcode}
architecture behaviour of parity_gen is
signal parity_s: std_logic;
begin
    parity <= parity_s;

    PARITY_PROCESS: process(bits) is
        variable acc: std_logic;
    begin
        acc := '0';
        for i in bits'range loop
            acc := acc xor bits(i);
        end loop;
        parity_s <= acc;
    end process;
end architecture;
\end{vhdlcode}

\noindent Now \code{acc := '0'} applies immediately; \code{i = 7} reads \code{'0'}, XORs in
\code{bits(7)} and updates \code{acc} at once; \code{i = 6} reads that just-updated value, and so on
through every iteration. The loop ends with \code{acc} holding the true XOR of all eight bits, and
\code{parity_s <= acc;} schedules that one correct value, the only assignment to \code{parity_s} in
the pass.

This is the on-disk worked example, \file{lectures/L05/parity_gen/parity_gen.vhd}, in full:

\vhdlfile{lectures/L05/parity_gen/parity_gen.vhd}

\subsection{What that loop is not}
\label{c5:loopnot}
It is not eight steps the hardware performs one after another. This is where a C background
misleads hardest, so it is worth stating outright: a \code{for} loop whose bounds are fixed at
elaboration is \textbf{unrolled} by the synthesis tool. Those eight iterations become eight XOR
gates in a chain, all settling within one propagation delay, in the same clock cycle. There is no
counter for \code{i}, no branch, no iteration in time. The loop writes eight gates in three lines;
it does not do eight things in sequence.

That is also why \code{acc}, which reads like a running total, costs no storage at all. It never has
to survive a clock edge: each of its values is just the wire between two of those XOR gates. A
\code{variable} is the right tool precisely because it never outlives one pass.

The same goes for \code{bits'range}. The \code{'range} \textbf{attribute} gives the index range of
\code{bits}, here \code{7 downto 0}, so the loop covers every bit without writing the bounds twice.
It is resolved at elaboration, like everything else in the loop header, which is exactly why the
tool can unroll it. \code{for i in 0 to 1 loop} in Chapter~\ref{c3:ch}'s \file{led_toggle.vhd}
means the same thing: both bits get their own gates and both update in the same cycle.

\subsection{The rule of thumb}
\label{c5:rule}
\begin{itemize}
  \item Use a \code{variable} for a scratch computation consumed entirely within one pass: loop
    accumulators, temporary swaps, intermediate values handed off to a signal at the end, exactly
    like \code{acc}.
  \item Use a \code{signal} for anything another process or the outside world reads, and anything
    that must hold its value across separate passes (a counter, or a shift register's stored bits,
    both in Chapter~\ref{c6:ch}).
\end{itemize}
A variable \emph{can} retain its value between passes through the same process, behaving like
static local storage rather than a fresh local each time, so it is possible to build state with
one. But because its updates are immediate rather than scheduled, it behaves differently from a
signal used the same way, particularly for anything modelling a chain of registers stepping one
stage per clock edge. If state genuinely needs to persist and be read by other logic, reach for a
\code{signal}: the scheduled-update semantics is what makes Chapter~\ref{c6:ch}'s counters and
shift registers work.

\section{What the FPGA Actually Builds}
\label{c5:a3}

\subsection{The two things the fabric is made of}
\label{c5:fabric}
An FPGA is not a blank slate onto which gates are etched. It is a fixed grid of identical small
blocks, manufactured before anybody wrote your VHDL, and ``synthesis'' is working out how to
configure them so they behave like the circuit you described. Two matter here:
\begin{itemize}
  \item A \textbf{look-up table}, or \textbf{LUT}: a tiny memory, typically four to six inputs and
    one output. Load it with the truth table of any Boolean function of those inputs and it computes
    that function. This is the honest answer to ``what does a gate become on an FPGA'': nothing
    becomes a gate. A four-input LUT computes \code{a and b}, \code{a xor b or (c and not d)}, or
    any other function of up to four inputs at exactly the same cost, because in every case it is
    the same block holding a different table.
  \item A \textbf{flip-flop}: one bit of clocked storage, sitting next to each LUT.
\end{itemize}
On the DE0-CV's Cyclone V these come packaged together as an \textbf{ALM}, holding a fracturable
six-input LUT, usable as one six-input function or two smaller ones, plus four flip-flops. That is
the unit the Quartus fitter report counts, which is why it talks about ALMs rather than gates.

That is nearly the whole story. \code{register4} becomes four flip-flops, \code{mux_8to1} a handful
of LUTs, and Chapter~\ref{c8:ch}'s state machine LUTs feeding the flip-flops that hold the state.

It also explains something from Chapter~\ref{c2:ch} that would otherwise look like a broken
promise: the book had you minimize expressions with Karnaugh maps to save gates, then never
mentioned the saving again. On an FPGA a four-input function costs one LUT whether you minimized it
or not. Minimization still matters: it is how you understand a circuit, it is exactly right for the
discrete-logic and ASIC flows K-maps were invented for, and reducing a function below the LUT's
input width is what stops it needing a second level of LUTs. But on this board it usually buys you
nothing, and it is better to know that than to wonder.

\subsection{Why a clock has a speed limit}
\label{c5:fmax}
Signals take time to get through a LUT and along the wires between blocks. That \textbf{propagation
delay} is why a clock cannot go arbitrarily fast.

Picture the shape every design in this book has, a flip-flop, some combinational logic, another
flip-flop:
\begin{itemize}
  \item On a rising edge the first presents a new value, which has to travel through every LUT and
    wire between the two and \emph{settle} before the next edge arrives.
  \item If it has not, the second flip-flop captures a value that is still changing, and the design
    does the wrong thing reliably and mysteriously.
\end{itemize}
So the clock period has to cover three things:
\begin{itemize}
  \item The first flip-flop's \textbf{clock-to-Q} delay: the time between the edge and its output
    actually changing.
  \item The \textbf{combinational delay} through the LUTs and wires between the two.
  \item The second flip-flop's \textbf{setup time}.
\end{itemize}
The slowest such path is the \textbf{critical path}, and one over that total is the \textbf{Fmax}.

Quartus computes Fmax for you, in the timing analyzer stage of the Quartus workflow
(Appendix~\ref{app:quartus}), but only if you tell it what the clock is. That means an \code{.sdc}
constraints file with a \code{create_clock} on the \code{50 MHz} input. Without one the analyzer has
no period of yours to compare anything against, so it warns that no base clocks were constrained
and reports Fmax against a default clock instead: a number, but not one you can act on. Separately,
a design that ``fails timing'' is one where some path turned out slower than the period you did
declare.

\textbf{None of that is part of this book.} You will not write an \code{.sdc}, read a timing
report, or meet a design that fails timing, because every design here fits inside 20\,ns with room
to spare. It is worth knowing the step exists and what it needs, so that you recognize it the first
time a design of yours gets close to the limit.

\noindent Two consequences worth carrying forward:
\begin{itemize}
  \item \textbf{Deep combinational logic costs you speed, not the number of flip-flops.} Chaining a
    long expression between two registers lengthens the critical path; splitting it so part of the
    work happens on one cycle and the rest on the next shortens it. That trade, more registers for a
    faster clock, is the single most common move in high-speed digital design.
  \item \textbf{\code{50 MHz} is a 20\,ns budget.} That is the number every path in your design is
    racing against, and it is a lot of time at these speeds.
\end{itemize}

\subsection{The latch you get by accident}
\label{c5:latch}
Chapter~\ref{c3:ch} built a D latch, showed it was transparent while enabled, then set it aside in
favour of the flip-flop. Here it comes back uninvited.

\Secref{c2:a5} required a \code{when others} on every \code{case}, and gave the language's reason: a
\code{case} has to cover every value of its selector's type. The general rule behind it is that a
combinational \code{process} must assign its output on \emph{every} path through it, and that
reason is now sayable:
\begin{itemize}
  \item VHDL requires a signal to keep its old value if nothing assigns a new one.
  \item If some path through your process leaves \code{x} unassigned, you have said \code{x} must
    \emph{remember} its previous value on that path, and remembering is not something combinational
    logic can do.
  \item So the tool builds the only thing that can: a latch, exactly the one from \secref{c3:a2}.
\end{itemize}
That is an \textbf{inferred latch}, and it is almost never what anybody meant. It is a warning
rather than an error, the design will often appear to work, and the message is worth recognizing:

\begin{console}
Warning: Inferred latch(es) for signal "x"
\end{console}

\noindent The fix is never to add a latch on purpose. It is to assign the signal on every path,
which is the rule you were already following.

\section{Looking Ahead}
\label{c5:a4}
Nothing in Chapters~\ref{c6:ch} to~\ref{c8:ch} introduces a new idiom. Counters, shift registers and
timers all hold state across clock edges using \secref{c3:a6}'s clocked process and the scheduling
rule above, and Chapter~\ref{c8:ch}'s state machines add only an enumerated type, a \code{case} on
it, and one new concurrent form, the conditional signal assignment. What changes is what the state
is used for.

Two things to watch for. Chapter~\ref{c6:ch}'s \code{serial_rx8} is the first design whose
correctness you cannot check by eye on an LED, which is where the provided testbenches stop being a
formality. And \secref{c4:a8}'s generic goes to work immediately: Chapter~\ref{c7:ch}'s \code{timer}
uses one for its tick count, and \code{walking_led} instantiates that timer alongside
\code{reset_sync} and \code{button_sync}, writing almost no logic of its own.
